\documentclass[10pt,a4paper]{amsart}
\usepackage[margin=19mm]{geometry}
\usepackage{amsmath,amssymb,mathtools}
\usepackage[scr=rsfs,cal=euler]{mathalfa}
\usepackage{xfrac}
\usepackage[unicode,hidelinks,bookmarks=false]{hyperref}
\newcommand{\ie}{i.e.\ }
\newcommand{\cf}{cf.\ }
\newcommand{\resp}{respectively\ }
\newcommand{\Subsection}{\subsection}
\providecommand{\nobreakdash}{\nobreak\hbox{-}}
\setlength{\emergencystretch}{2em}
\multlinegap0pt
\usepackage{bm}
\usepackage{bbm}
\usepackage{dsfont}
\usepackage{xspace}
\usepackage{cancel}
\usepackage{mathtools}
\usepackage{amsthm}
\makeatletter
\@ifundefined{theorem}{\newtheorem{theorem}{Theorem}}{}
\@ifundefined{claim}{\newtheorem{claim}[theorem]{Claim}}{}
\@ifundefined{lemma}{\newtheorem{lemma}[theorem]{Lemma}}{}
\@ifundefined{proposition}{\newtheorem{proposition}[theorem]{Proposition}}{}
\@ifundefined{remark}{\newtheorem{remark}[theorem]{Remark}}{}
\makeatother
\DeclareMathOperator{\vol}{vol}
\newcommand{\Gr}[2]{G(#1,#2)}        % Grassmannian G(j,n)
\newcommand{\R}{\mathbb{R}}
\newcommand{\Sph}{\mathbb{S}}        % unit sphere; use \Sph^{n-1}
\newcommand{\cF}{\mathcal{F}}
\newcommand{\cH}{\mathcal{H}}        % homogeneous polynomials
\newcommand{\cP}{\mathcal{P}}        % positive forms
\newcommand{\opV}{\mathds{V}}       % operator in section 4
\title[Computing intrinsic volumes of sublevel sets and application]{Computing intrinsic volumes of sublevel sets and applications}
\author{Tr\'i Minh L\^e, Khai-Hoan Nguyen-Dang}
\date{Source: arXiv:2510.24001v2 (CC BY 4.0)}
\begin{document}
\maketitle

\noindent\emph{Source and licence.} Computing intrinsic volumes of sublevel sets and applications. Version arXiv:2510.24001v2 (CC BY 4.0), distributed under the Creative Commons Attribution 4.0 International licence, as recorded in the arXiv licence metadata for this exact version and shown on the abstract page https://arxiv.org/abs/2510.24001v2. Full rights notice: licenses/arxiv-2510.24001-CC-BY-4.0.txt.\par\smallskip

\noindent\textit{Editorial scope.} Counted unit: the Lemma whose source label is \texttt{lem.Vtil-gateaux} in the article (source0.tex lines 1209--1219, source label \texttt{lem.Vtil-gateaux}), delivered together with the article's own complete original proof of it (source0.tex lines 1221--1271, one \texttt{proof} environment of 51 lines). No mathematics is written, repaired or re-derived: statement and proof are the article's own text line for line. Required context retained uncounted: rem:integrability (lines 606--642); prop.cF-coer (lines 2261--2270). Every definition, display and label of those lines is retained unchanged. Omitted from this package and not redistributed: 1--605, 643--1208, 1220--1220, 1272--2260, 2271--2370. Nothing inside a retained excerpt is omitted or altered: every retained body is delivered exactly as the article's lines are written, so no omission receipt of any kind accompanies this package. Citations: the retained excerpts carry no citation command at all, so no bibliography entry of the article is transcribed and no reference is redistributed. Reference adaptation: none was needed and none was made; every reference inside the delivered unit resolves inside it, so no unresolved reference or citation is produced. Printed slips kept literal: no printed word, symbol, hypothesis or number of the counted statement or of its proof was altered. Layout and class: the article's own class is \texttt{article} with packages including amsmath, amsfonts, amssymb, mathtools, amsthm, xcolor, float, enumitem, cleveref, euscript, resmes, dsfont, mathrsfs, fontenc; the wrapper is the lane's pinned \texttt{amsart} editorial wrapper at 10pt on A4, which restates the theorem-like environments the retained text needs and adds the article's own macro definitions that the retained text needs (\texttt{\textbackslash Gr}, \texttt{\textbackslash R}, \texttt{\textbackslash Sph}, \texttt{\textbackslash cF}, \texttt{\textbackslash cH}, \texttt{\textbackslash cP}, \texttt{\textbackslash opV}, \texttt{\textbackslash vol}), with extra wrapper packages (bm, bbm, dsfont, xspace, cancel, mathtools). The delivered numbering is the unit's own. Assets: no retained span contains a figure environment, a graphics-inclusion command, a PDF-page inclusion, a table or a TikZ picture; no image file is copied into this package, the package declares no asset dependency and no image file is redistributed with it. Licence: version arXiv:2510.24001v2 of this article is distributed under the Creative Commons Attribution 4.0 International licence, as recorded in the arXiv licence metadata for this exact version and shown on the abstract page https://arxiv.org/abs/2510.24001v2. A full rights notice is delivered at licenses/arxiv-2510.24001-CC-BY-4.0.txt. Characters: every retained excerpt is pure ASCII, so the delivered engine, XeLaTeX with its default Latin Modern text font, reports no missing character.
\begin{remark}[Integrability]\label{rem:integrability}
\textit{(i)} For any $f\in \cF_{n,d}$ and any $E\in G(j,n)$, it holds
\[
  \int_E \exp(-\mathscr R_E f(y))\,dy < + \infty\quad \text{ and } \quad 
  \int_E \exp(-\Pi_E f(y))\,dy< +\infty.
\]
Indeed, thanks to Proposition~\ref{prop.cF-coer}, for any fixed $f \in \cF_{n, d}$, there exists $\varpi > 0$ such that $f(x) \geq \varpi \| x \|^d$ for every $x \in \R^n$.
Hence, we get
\begin{align*}
    \int_E \exp(- \mathscr R_E f(y)) \, dy \leq \int_E \exp( - \varpi \| y \|^d) \, dy = \frac{j }{d} \Gamma(j/d) \vol_j(B^j) \varpi^{-j/d} < + \infty
\end{align*}

Furthermore, with a direct computation, we have
\begin{equation*}
    \| y + z \|^2 = \| y \|^2 + 2 \underbrace{\langle y, z \rangle}_{\, = \, 0} + \| z \|^2 \geq \| y \|^2, \quad \text{ for every  } y \in E, \, z \in E^\perp,
\end{equation*}
which leads to $\textstyle\inf_{z \in E^\perp} \| y + z \|^d = \| y \|^d$ for every $y \in E$. 
Remark that this fact may be false if we consider other norms instead of Euclidean norm.
Therefore, we have
\[
    \Pi_E f(y) = \inf_{z \in E^\perp} f(y + z) \geq \varpi \inf_{z \in E^\perp} \| y + z \|^d = \varpi \| y \|^d,
\]
which leads to
\begin{align*}
    \int_E \exp( - \Pi_E f(y)) \, dy \leq \int_E \exp(- \varpi \| y \|^d) \, dy < +\infty.
\end{align*}

\medskip

\textit{(ii)}
As a consequence of \textit{(i)}, since $\nu_j$ is a probability measure on $G(j,n)$, we also have, for every $f \in \cF_{n, d}$,
\begin{align*}
	& \int_{G(j,n)} \ \int_{E} \exp\big(-\Pi_E f(y)\big)\,dy \, d\nu_j(E) < + \infty \\
    \quad \text{ and }   & \int_{G(j,n)} \ \int_{E} \exp\big(-\mathscr R_E f(y)\big)\,dy \, d\nu_j(E) < + \infty.
\end{align*}
%by Lemma~\ref{lem.vol_j} and Proposition~\ref{prop.pro-sec}; in particular all integrals below are finite.
\end{remark}

\begin{proposition}\label{prop.cF-coer}
Let $f \not\equiv 0$ be a lower semicontinuous nonnegative and positively $d$--homogeneous function.
Then, the following assertions are equivalent:
\begin{itemize}
	\item[(i)] $[f \leq 1]$ is bounded (equivalently, compact);
	\item[(ii)] $\textstyle\min_{e \in \Sph^{n - 1}} f(e) > 0$;
	\item[(iii)] there exists $\varpi > 0$ such that $f(x) \geq \varpi \| x \|^d$, for every $x \in \R^n$;
	\item[(iv)] (definition of $\cP_{n, d}$) $f(x) > 0$ for every $x \neq 0$.
\end{itemize}
\end{proposition}

\begin{lemma}[G\^ateaux derivative of $\widetilde\opV_j$]\label{lem.Vtil-gateaux}
    Let $f \in \cP_{n, d}$ and let $1 \leq j \leq n - 1$.
    Then, $\widetilde\opV_j$ is G\^ateaux differentiable and moreover for any $\phi \in \cH_{n, d}$ one has
    \begin{align*}
       \mathrm d \widetilde\opV_j(f; \phi) = - \dfrac{\sigma_{j, n}}{\Gamma(1 + j/d)} \int_{G(j, n)}\int_E (\mathscr{R}_E \phi)(y) \exp(- \mathscr R_E f(y)) \, dy \, d\nu_j(E).
    \end{align*}
%    In particular, if $\{ \phi_k \}_{ k \in I}$ is the canonical basis of $\cH_{n, d}$, then
%    \begin{align*}
%    		\dfrac{\partial \widetilde\opV_j}{\partial \phi_k}(f) := \mathrm d \widetilde\opV_j(f; \phi_k) = - \dfrac{\sigma_{j, n}}{\Gamma(1 + j/d)} \int_{G(j, n)}\int_E (\mathscr{R}_E \phi_k)(x) \exp(- \mathscr R_E f(x)) \, dx \, d\nu_j(E).
%    \end{align*}
\end{lemma}

\begin{proof}
The case $\phi \equiv 0$ is vacuous, so we assume $\phi \not\equiv 0$. 
To proceed, we need a lower bound for the operator $\mathscr{R}_E f$ under small perturbations.

\begin{claim}\label{claim.gateaux-bound}
There exists $\bar t > 0$ independent of $E$ such that
\begin{align*}
	\mathscr R_E (f + t \phi)(y) \geq \frac{\varpi}{2} \| y \|^d, \quad \text{ for every } y \in E, \, |t| \leq \bar t.
\end{align*}
\end{claim}

\textit{Proof of Claim~\ref{claim.gateaux-bound}.}
Thanks to Remark~\ref{rem:integrability}, there exists $\varpi > 0$ independent of $E$ such that $\mathscr R_E f(y) \geq \varpi \| y \|^d$ for every $y \in E$ and $E \in \Gr{j}{n}$.
For a nonzero $\phi \in \cH_{n, d}$, set $M = \textstyle \| \phi \|_{L^\infty(\mathbb S^{n - 1})} \in (0, + \infty)$.
Choose $\bar t = \varpi/2M$.
Thanks to the $d$--homogeneity of $\phi$, for any $|t| \leq \bar t$, we have
\[
	|t| |\phi(y)| = |t| \| y \|^d |\phi(y / \| y \|)| \leq M \bar t \| y \|^d, \quad \text{ for every } y \in E \setminus \{ 0 \}.
\]
Then, for any $|t| \leq \bar t$, a direct computation leads to
\begin{align*}
	\mathscr R_E (f + t \phi)(y) = \mathscr R_E f(y) + t \mathscr R_E \phi(y) \geq \varpi \| y \|^d - |t| |\phi(y)| \geq \frac{\varpi}{2} \| y \|^d, \quad \text{ for every } y \in E,
\end{align*}
which completes the proof. \hfill  $\lozenge$

On the one hand, note that
\[
	\dfrac{d}{dt} \exp(- \mathscr R_E(f + t \phi)(y)) = - \mathscr R_E \phi(y) \exp( - \mathscr R_E(f + t \phi)(y)), \quad \text{ for every } y \in  E.
\]
On the other hand, the $d$--homogeneity of $\phi$ implies $| \mathscr R_E \phi(y) | \leq M \| y \|^d$.
Hence, combining with Claim~\ref{claim.gateaux-bound}, we have
\begin{align*}
	\left\vert (\mathscr R_E\phi)(y)\,\exp(-\mathscr R_E (f + t \phi)(y)) \right\vert
\ \le\ M(1+\|y\|^d)\,e^{-(\varpi/2)\|y\|^d}, \quad \text{ for every } |t| \leq \bar t, \, y \in E.
\end{align*}
A direct computation leads to
\[
    \int_{G(j, n)} \int_E (1 + \| y \|^d) e^{-(\varpi/2)\|y\|^d} \, dy \, d\nu_j(E) < + \infty.
\]
Therefore, by the dominated convergence theorem, we get
\begin{align*}
       \mathrm d \widetilde\opV_j(f; \phi) = - \dfrac{\sigma_{j, n}}{\Gamma(1 + j/d)} \int_{G(j, n)}\int_E (\mathscr{R}_E \phi)(y) \exp(- \mathscr R_E f(y)) \, dy \, d\nu_j(E).
\end{align*}
Note that the linearity of $\phi \mapsto \mathrm d \widetilde\opV_j(f; \phi)$ follows from the linearity of the operator $\mathscr R_E$.
Moreover, one can directly check that
\begin{align*}
	|  \mathrm d \widetilde\opV_j(f; \phi)| \leq \dfrac{\sigma_{j, n}}{\Gamma(1 + j/d)}\| \phi \|_{L^\infty(\Sph^{n - 1})} \underbrace{\int_{\R^j} \| x \|^d \exp(- \varpi \| x \|^d) dx}_{\,  < \, + \infty},
\end{align*}
which leads to the boundedness of the map $\phi \mapsto \mathrm d \widetilde\opV_j(f; \phi)$.
Lastly, we can conclude that $\widetilde\opV_j$ is G\^ateaux differentiable, which finishes the proof of Lemma~\ref{lem.Vtil-gateaux}.
\end{proof}

\end{document}
